\section{Introduction}
\label{sec:introduction}

% Superconducting qubits require frequent calibration because transition
% frequencies drift with flux bias, temperature, control conditions, and the
% electromagnetic environment.  Frequency error appears directly as a detuning
% in single-qubit control and propagates into gate, readout, and higher-level
% calibration tasks~\cite{krantz2019quantum,blais2021circuit}.  Modern calibration
% systems therefore organize repeated measurements and parameter updates into
% automated feedback loops~\cite{kelly2018physical,wittler_integrated_2021}.
% Reducing the cost of the elementary frequency measurement is valuable because
% the same primitive is executed across devices and repeated as the processor
% drifts~\cite{rol2017restless,werninghaus_high-speed_2021}.  This motivates a
% broader question: can the short control-pulse sequences already used around a
% working point provide calibration information themselves?

% Ramsey spectroscopy is a robust frequency estimator.  A pair of $\pi/2$
% pulses separated by a variable free-evolution interval converts detuning into
% population oscillations, and the oscillation frequency is extracted from a
% scan.  The scan provides a useful capture range and high precision, but its
% cost scales with the number of delay points and, when two artificial detunings
% are used to determine the sign, with two complete scans.  Fixed-delay Ramsey
% measurements reduce this burden once the frequency is already known locally,
% and have been incorporated into closed-loop stabilization protocols
% ~\cite{vepsalainen2022improving}.  Their local phase response, however, becomes
% ambiguous after phase wrapping.

% Here we test this possibility in a frequency-calibration setting.  We remove
% the free-evolution interval and infer detuning from the response accumulated
% during two finite control pulses.  Orthogonal analysis branches form an odd
% differential signal whose slope is determined by a response kernel.  Such
% control-time response functions connect temporal sensing to driven-qubit
% dynamics~\cite{herb2024quantum,lin2025application}.  They suggest that a short
% pulse can be more than an operation surrounding a calibration experiment: its
% transient response can itself be the probe.  The question is then not only
% whether this response contains frequency information, but whether that local
% information can remain valid as a calibration loop moves the device.

% Two constraints are central.  First, the short-pulse response is nonlinear and
% only locally invertible; its leading correction depends on the full three-time
% response kernel.  Second, if the feedback loop changes the flux bias while the
% drive remains fixed, the qubit can leave the discriminator's monotonic window,
% corrupting both the frequency estimate and the subsequent controller update.

% Our central result is that short-pulse sequences can participate in
% closed-loop frequency calibration when these two constraints are controlled.
% A full three-time response kernel supplies the local cubic inverse.  A secant
% estimate of the frequency--flux slope then predicts the transition frequency
% at the next bias point.  Driving at this predicted frequency keeps the pulse
% response inside its calibrated window, while the feedback residual remains
% defined relative to the fixed target.  The response kernel and drive tracking
% therefore play supporting but distinct roles: one extracts information from
% the short sequence, and the other preserves the validity of that information
% along the calibration trajectory.

% We first derive the short-pulse discriminator and its local inverse.  We then
% separate drive tracking, which preserves measurement validity, from the flux
% update, which closes the control loop.  Numerical tests compare the local
% probe with Ramsey strategies using frequency errors recomputed from the
% analytic simulation model and instrumented solver counts.  This ground-truth
% check is separate from the independent measurement protocol specified for
% \textsc{Verify}. Besides demonstrating a useful operating regime,
% the comparison exposes a negative result: simply placing a short-pulse method
% before Ramsey in a coarse-to-fine workflow can cost more than Ramsey alone.
INTRODUCTION PLACEHOLDER: The introduction will be finalized after the main text is confirmed.  It will motivate the work, summarize the method, and highlight the key contributions and implications for superconducting-qubit frequency calibration.

\newpage
